\documentclass[11pt]{article}
\usepackage[margin=1in]{geometry}
\usepackage{hyperref}
\usepackage{enumitem}
\usepackage{xcolor}
\usepackage{tikz}
\usetikzlibrary{shapes.geometric, arrows.meta, positioning, shadows}

\tikzset{
  startstop/.style={ellipse, draw=teal!70!black, thick, fill=teal!15,
    minimum width=2.4cm, minimum height=0.9cm, align=center,
    font=\small\bfseries, drop shadow},
  process/.style={rectangle, rounded corners=3pt, draw=blue!50!black, thick,
    fill=blue!6, text width=6.6cm, minimum height=1.05cm, align=center,
    font=\small, inner sep=6pt, drop shadow},
  io/.style={trapezium, trapezium left angle=68, trapezium right angle=112,
    draw=orange!70!black, thick, fill=orange!10, text width=6.6cm,
    minimum height=1.0cm, align=center, font=\small, inner sep=6pt, drop shadow},
  arrowstyle/.style={-{Latex[length=2.3mm]}, thick, gray!55!black},
}

% #1 init text, #2 save text, #3 rk4 text, #4 clip/adjust text, #5 end text
\newcommand{\DnsFlow}[5]{%
\begin{center}
\begin{tikzpicture}[node distance=6mm and 6mm]
  \node[startstop] (start) {Start};
  \node[process, below=of start] (init) {#1};
  \node[io, below=of init] (save) {#2};
  \node[process, below=of save] (rk4) {#3};
  \node[process, below=of rk4] (clip) {#4};
  \node[process, below=of clip] (enforce) {Restore the divergence-free condition on $u$};
  \node[startstop, below=of enforce] (endnode) {#5};
  \draw[arrowstyle] (start) -- (init);
  \draw[arrowstyle] (init) -- (save);
  \draw[arrowstyle] (save) -- (rk4);
  \draw[arrowstyle] (rk4) -- (clip);
  \draw[arrowstyle] (clip) -- (enforce);
  \draw[arrowstyle] (enforce) -- (endnode);
  \draw[arrowstyle] (enforce.east) -- ++(1.5,0) |- (save.east)
    node[pos=0.27, right, font=\scriptsize, align=left, text=gray!50!black]
    {repeat until\\[-2pt]$t \ge T$};
\end{tikzpicture}
\end{center}
}

\title{File Types in the Turbulence--Droplet DNS Code}
\author{}
\date{}

\begin{document}
\maketitle

\section{Overview}
This code simulates forced, statistically homogeneous, isotropic
turbulence with a Fourier (pseudo-spectral) method, one-way coupled
to particles that mimic cloud droplets. The velocity field is evolved
in Fourier space; particles are represented either as a continuum
number-density field advected on the ``slow manifold'' (small Stokes
number, i.e.\ particles that closely follow the fluid) or as explicit
tracked particles (finite Stokes number, referred to as ``big''
particles). The code runs in parallel across many processors using
MPI, with the simulation domain sliced into 1-D slabs along $x$.

\section{Main simulation scripts (\texttt{forced-dns-*.py})}
These are five independent, largely duplicated top-level scripts,
not a shared library with switches. Each hardcodes its own grid
size, time step, save path, and physics; each is a program you run
directly with \texttt{mpirun}, not a module you import.

\begin{itemize}[leftmargin=*]
  \item \textbf{\texttt{forced-dns-sm.py}} --- the baseline: forced
    Navier--Stokes turbulence plus the slow-manifold number-density
    field \texttt{n}, with no explicit particles. This is the file
    whose style the others imitate.
  \item \textbf{\texttt{forced-dns-sm-clip-new.py}} --- identical to
    the baseline except that the crude step of clipping negative
    \texttt{n} to zero and rescaling is replaced by a root-finding
    procedure (\texttt{clip\_zero}) that clips to an exact target
    mean.
  \item \textbf{\texttt{forced-dns-log-sm.py}} --- evolves
    $\log(n)$ instead of \texttt{n} directly, which avoids needing
    to clip for positivity at all (a logarithm is automatically
    positive when exponentiated). Uses a different time step,
    save interval, and forcing amplitude, and its evolution equation
    carries different terms.
  \item \textbf{\texttt{forced-dns-sm-big.py}} --- adds explicit
    tracked (``big'') particles for several Stokes numbers and
    initial conditions simultaneously, in addition to the
    slow-manifold field. This is the only variant that places new
    particles using the $Q$-criterion (a way of seeding particles in
    regions of strong flow rotation).
  \item \textbf{\texttt{forced-dns-sm-big-rndm.py}} --- a big-particle
    variant with added random forcing on the particles themselves.
    It carries no density field at all.
\end{itemize}

Each script follows the same internal layout: parameter block and
grid/Fourier-transform setup, a forcing routine, a clipping routine,
the spectral right-hand side of the evolution equations, a
fourth-order Runge--Kutta time-stepper, a checkpoint-saving routine,
and a main time-stepping loop, followed by construction of the
initial conditions and the call that runs the simulation.

\subsection{Control flow of each variant}
All five scripts share the same outer skeleton: build or load the
starting state, then repeatedly take a Runge--Kutta step, fix up the
result, and save on a schedule, until the end time is reached. The
flowcharts below highlight exactly where each variant departs from
that skeleton.

\paragraph{\texttt{forced-dns-sm.py}}
\DnsFlow{Load a restart, or build a fresh field ($u$, $n$)}%
{Every \texttt{st} steps: save $u$, $n$, and the energy/flux spectra}%
{RK4 step: forced Navier--Stokes for $u$; advect $n$ with the flow,
corrected for the particles' inertia}%
{Clip negative $n$ to zero, rescale so its mean is unchanged
(\texttt{clip\_zero})}%
{$t \ge T$: save the final state and stop}

\paragraph{\texttt{forced-dns-sm-clip-new.py}}
\DnsFlow{Load a restart, or build a fresh field ($u$, $n$)}%
{Every \texttt{st} steps: save $u$, $n$, and the energy/flux spectra}%
{RK4 step: forced Navier--Stokes for $u$; advect $n$ with the flow,
corrected for the particles' inertia}%
{Root-find a clip of $n$ that hits a target mean exactly
(\texttt{clip\_zero})}%
{$t \ge T$: save the final state and stop}

\paragraph{\texttt{forced-dns-log-sm.py}}
\DnsFlow{Load a restart, or build a fresh field ($u$, $\log n$)}%
{Every \texttt{st} steps: save $u$, $\log n$, and the energy/flux spectra}%
{RK4 step: forced Navier--Stokes for $u$; advance $\log n$ using the
flow's velocity gradient and its divergence}%
{Clamp $\log n$ between a floor and a ceiling, rescale so the mean
of $n$ stays 1 (\texttt{clip\_error})}%
{$t \ge T$: save the final state and stop}

\paragraph{\texttt{forced-dns-sm-big.py}}
\DnsFlow{Load a restart, or place fresh particles by the $Q$-criterion,
for each Stokes number, alongside the field ($u$, $n$)}%
{Near each save time: write $u$, $n$, and every Stokes number's
particle state}%
{RK4 step: forced Navier--Stokes for $u$; for each Stokes number,
advect $n$ \emph{and} push its particles (interpolate the flow onto
them, integrate their motion)}%
{Clip negative $n$ to zero per Stokes number; shrink the time step to
resolve the fastest particle}%
{$t \ge T$: save the final state and stop}

\paragraph{\texttt{forced-dns-sm-big-rndm.py}}
\DnsFlow{Load a restart, or place fresh particles from the local flow
velocity, for each Stokes number (no density field)}%
{Near each save time: write $u$ and every Stokes number's particle
state}%
{RK4 step: forced Navier--Stokes for $u$; for each Stokes number,
push particles with the inertial particle equation}%
{Add one random velocity kick to every particle}%
{$t \ge T$: save the final state and stop}

\section{Shared support modules}
\begin{itemize}[leftmargin=*]
  \item \textbf{\texttt{initial\_conditions.py}} --- shared startup
    logic used by all five scripts. It either loads a previously
    saved state (a restart) or builds a fresh one, and is
    parameterized (via constructor arguments) to handle every
    difference between the five scripts' restart conventions:
    which checkpoint folder to use, how to clip the density, whether
    the field on disk is \texttt{n} or $\log n$, how particles are
    seeded from scratch, and so on.
  \item \textbf{\texttt{particles.py}} --- the tracked-particle
    engine (class \texttt{MPI\_particles}). Holds each particle's
    position, velocity, and mass; interpolates the fluid velocity
    onto particle positions and deposits particle quantities back
    onto the grid; exchanges particles between processors when they
    cross slab boundaries; and computes the particle equations of
    motion, including the random-forcing option. A few core
    interpolation routines are compiled just-in-time for speed and
    must stay compatible with that compiler's restrictions.
  \item \textbf{\texttt{restart.py}} --- an earlier, now-superseded
    version of the restart logic. No script imports it, so it has no
    effect on any run; it is kept only as a leftover.
\end{itemize}

\section{Analysis and plotting scripts}
\texttt{plot.py}, \texttt{plot-correlation.py}, and the dated
\texttt{plot-YYYYMMDD.py} files are standalone, notebook-style
scripts (organized in cells) that read simulation output and produce
figures. They are not imported by any other file and are treated as
disposable: a new plotting need is normally addressed by adding a
new dated file rather than editing an old one.

\section{Data}
\texttt{data\_cosine/} holds a small local sample of simulation
output, organized by forcing flag, grid resolution, and Reynolds
number, then by save time. Each save-time folder contains
Fourier-space velocity field pieces (one file per processor slab),
energy- and flux-spectrum diagnostics, and, per Stokes number,
either number-density files or tracked-particle state files. This
sample data is excluded from version control and is much smaller
than the full dataset, which lives on the cluster where the
simulations actually run.

\texttt{comparing\_spectra.hdf5} is a single self-describing binary
data file (HDF5 format) at the repository root, used for comparing
energy spectra.

\section{Repository metadata}
\begin{itemize}[leftmargin=*]
  \item \textbf{\texttt{CLAUDE.md}} --- instructions and
    architectural notes for the Claude assistant working in this
    repository.
  \item \textbf{\texttt{.gitignore}} --- lists paths excluded from
    version control (the sample data folder, compiled-code caches,
    and the ignore file itself).
  \item \textbf{\texttt{.claude/}} --- assistant configuration,
    including the agent that reviews each response's changes for
    unnecessary complexity.
\end{itemize}

\end{document}
